% ── §2.2: Ομοσπονδιακή Μάθηση (Federated Learning) ──────────

\section{Ομοσπονδιακή Μάθηση (Federated Learning)}
\label{sec:federated_learning}

Η Ομοσπονδιακή Μάθηση (Federated Learning, FL) είναι ένα παράδειγμα
κατανεμημένης μηχανικής μάθησης στο οποίο ένα καθολικό μοντέλο
εκπαιδεύεται από ένα πλήθος $n$ συμμετεχόντων (clients) χωρίς τα
ακατέργαστα δεδομένα τους να εγκαταλείπουν ποτέ τη συσκευή τους
\cite{McMahan2017}. Αντί για δεδομένα, μεταδίδονται μόνο ενημερώσεις
μοντέλου (model updates) στον κεντρικό συντονιστή (server), ο οποίος
τις συνθέτει σε ένα βελτιωμένο καθολικό μοντέλο.

\subsection{Αρχιτεκτονική Server-Client}
\label{subsec:fl_architecture}

Το τυπικό FL σύστημα αποτελείται από έναν \textbf{server} και $n$
\textbf{clients}. Ο server αρχικοποιεί ένα καθολικό μοντέλο
$\mathbf{w}^{(0)}$ και ορίζει τη στρατηγική επιλογής clients ανά γύρο
επικοινωνίας. Κάθε client $i$ διατηρεί τοπικό σύνολο δεδομένων
$\mathcal{D}_i$ με $n_i = |\mathcal{D}_i|$ δείγματα. Η κατανεμημένη
βελτιστοποίηση στοχεύει στην ελαχιστοποίηση της σταθμισμένης συνολικής
απώλειας:

\begin{equation}
    \label{eq:fl_objective}
    \min_{\mathbf{w}} \; F(\mathbf{w}) = \sum_{i=1}^{n}
    \frac{n_i}{N} \, F_i(\mathbf{w}), \quad N = \sum_{i=1}^{n} n_i
\end{equation}

\noindent όπου $F_i(\mathbf{w}) = \frac{1}{n_i}\sum_{(\mathbf{x},y)
\in \mathcal{D}_i} \ell(\mathbf{w}; \mathbf{x}, y)$ είναι η τοπική
απώλεια του client $i$ και $\ell(\cdot)$ η συνάρτηση κόστους
(cross-entropy για ταξινόμηση).

\subsection{Ο Αλγόριθμος FedAvg}
\label{subsec:fedavg}

Ο αλγόριθμος Federated Averaging (FedAvg), που εισήχθη από τους McMahan
et al.\ \cite{McMahan2017}, αποτελεί τον de facto αλγόριθμο αναφοράς
του FL. Επαναλαμβάνει τα ακόλουθα βήματα για $R$ γύρους επικοινωνίας:

\begin{enumerate}
    \item \textbf{Κοινοποίηση:} Ο server αποστέλλει τις τρέχουσες
    παραμέτρους $\mathbf{w}^{(r)}$ στο υποσύνολο $\mathcal{S}^{(r)}$
    επιλεγμένων clients (με $|\mathcal{S}^{(r)}| = m \leq n$).

    \item \textbf{Τοπική εκπαίδευση:} Κάθε client $i \in \mathcal{S}^{(r)}$
    εκτελεί $E$ εποχές τοπικής βελτιστοποίησης (SGD ή Adam) στα
    δεδομένα του, παράγοντας ενημερωμένες παραμέτρους
    $\mathbf{w}_i^{(r+1)}$.

    \item \textbf{Συνάθροιση:} Ο server συνθέτει τις ενημερώσεις
    στάθμισε με βάση το μέγεθος των τοπικών συνόλων:
    \begin{equation}
        \label{eq:fedavg}
        \mathbf{w}^{(r+1)} = \sum_{i \in \mathcal{S}^{(r)}}
        \frac{n_i}{\sum_{j \in \mathcal{S}^{(r)}} n_j}
        \mathbf{w}_i^{(r+1)}
    \end{equation}
\end{enumerate}

Η θεωρητική ανάλυση σύγκλισης του FedAvg υπό Non-IID κατανομές
παρουσιάζεται από τους Li et al.\ \cite{Li2020}, οι οποίοι αποδεικνύουν
ότι η ετερογένεια δεδομένων μπορεί να επιβραδύνει σημαντικά τη σύγκλιση
λόγω του φαινομένου \textit{client drift}. Το Σχήμα~\ref{fig:fl_round}
απεικονίζει σχηματικά τον κύκλο ενός γύρου FL.

\begin{figure}[H]
\centering
\begin{tikzpicture}[
    node distance=1.5cm and 2.2cm,
    box/.style={rectangle, rounded corners=4pt, draw=black, thick,
                minimum width=2.8cm, minimum height=0.9cm,
                align=center, font=\small},
    server/.style={box, fill=blue!15},
    client/.style={box, fill=orange!20},
    arrow/.style={-{Stealth[length=6pt]}, thick},
    dasharrow/.style={-{Stealth[length=6pt]}, thick, dashed}
]

% Server
\node[server] (server) {FL Server\\[2pt]$\mathbf{w}^{(r)}$};

% Clients
\node[client, below left=1.8cm and 2.5cm of server]  (c1) {Client 1\\[2pt]$\mathcal{D}_1$};
\node[client, below=1.8cm of server]                  (c2) {Client 2\\[2pt]$\mathcal{D}_2$};
\node[client, below right=1.8cm and 2.5cm of server]  (c3) {Client $N$\\[2pt]$\mathcal{D}_N$};

% Dots between clients
\node[font=\Large] at ($(c2)!0.5!(c3) + (0.6,0)$) {$\cdots$};

% Broadcast arrows (server → clients)
\draw[arrow, blue!60] (server) -- node[left,  font=\scriptsize]{$\mathbf{w}^{(r)}$} (c1);
\draw[arrow, blue!60] (server) -- node[right, font=\scriptsize]{$\mathbf{w}^{(r)}$} (c2);
\draw[arrow, blue!60] (server) -- node[right, font=\scriptsize]{$\mathbf{w}^{(r)}$} (c3);

% Local training labels
\node[font=\scriptsize, text=gray] at ($(c1)+(0,-0.7)$) {$E$ τοπικές εποχές};
\node[font=\scriptsize, text=gray] at ($(c2)+(0,-0.7)$) {$E$ τοπικές εποχές};
\node[font=\scriptsize, text=gray] at ($(c3)+(0,-0.7)$) {$E$ τοπικές εποχές};

% Aggregation node
\node[server, below=3.8cm of server] (agg) {FedAvg\\[2pt]Aggregation};

% Upload arrows (clients → aggregation)
\draw[arrow, orange!70!black] (c1) -- node[left,  font=\scriptsize]{$\mathbf{w}_1^{(r+1)}$} (agg);
\draw[arrow, orange!70!black] (c2) -- node[right, font=\scriptsize]{$\mathbf{w}_2^{(r+1)}$} (agg);
\draw[arrow, orange!70!black] (c3) -- node[right, font=\scriptsize]{$\mathbf{w}_N^{(r+1)}$} (agg);

% Loop arrow
\draw[dasharrow, green!50!black, bend right=50]
      (agg.east) to node[right, align=left, font=\scriptsize]{$\mathbf{w}^{(r+1)}$\\[1pt]επόμενος γύρος} (server.east);

\end{tikzpicture}
\caption{Κύκλος ενός FL γύρου (FedAvg): ο server εκπέμπει τις τρέχουσες παραμέτρους, κάθε client εκτελεί $E$ τοπικές εποχές χωρίς κοινοποίηση δεδομένων, και τα ενημερωμένα βάρη συνδυάζονται σταθμισμένα για το επόμενο global μοντέλο.}
\label{fig:fl_round}
\end{figure}

\subsection{Προκλήσεις του Federated Learning}
\label{subsec:fl_challenges}

Η ολοκληρωμένη ανασκόπηση των Kairouz et al.\ \cite{Kairouz2021}
κατατάσσει τις κύριες προκλήσεις του FL σε τρεις κατηγορίες:

\begin{itemize}
    \item \textbf{Ετερογένεια δεδομένων (Non-IID):} Τα δεδομένα κάθε
    client ακολουθούν διαφορετική κατανομή, καθώς αντικατοπτρίζουν
    τις ατομικές συνήθειες και κινητικά μοτίβα κάθε χρήστη. Αυτό
    αποκλίνει από τη βέλτιστη υπόθεση IID (Independent and Identically
    Distributed) της στατιστικής θεωρίας \cite{Li2020}.

    \item \textbf{Κόστος επικοινωνίας:} Κάθε γύρος απαιτεί μετάδοση
    ολόκληρου του μοντέλου (ή των παραμέτρων του) μεταξύ server και
    clients. Για μεγάλα μοντέλα ή αργές συνδέσεις, αυτό αποτελεί
    σημαντικό περιορισμό \cite{Konecny2016}.

    \item \textbf{Αξιοπιστία clients:} Clients μπορεί να αποσυνδεθούν
    κατά τη διάρκεια ενός γύρου (stragglers) ή να υποβάλουν ψευδείς
    ενημερώσεις (Byzantine clients). Η παρούσα εργασία αντιμετωπίζει
    την πρώτη κατηγορία μέσω του μηχανισμού κινήτρων
    (Κεφάλαιο~\ref{ch:implementation}).
\end{itemize}

\subsection{Ιδιωτικότητα στο FL}
\label{subsec:fl_privacy}

Παρόλο που το FL δεν μεταδίδει ακατέργαστα δεδομένα, οι ενημερώσεις
μοντέλου μπορούν υπό συνθήκες να αποκαλύψουν ευαίσθητες πληροφορίες
μέσω \textit{gradient inversion attacks} \cite{Shokri2015}. Μέθοδοι
όπως η Differential Privacy \cite{Abadi2016} και η ασφαλής συνάθροιση
\cite{Bonawitz2017} αμβλύνουν αυτόν τον κίνδυνο, αν και βρίσκονται
εκτός του άμεσου πεδίου της παρούσας εργασίας.
